\chapter{Engenharia de Software e Modernização}
\section{Engenharia de Requisitos com Sommerville}
O \verb|app.py| preserva integralmente as regras de negócio do legado Vaadin/Java, seguindo as quatro fases clássicas propostas por Sommerville~\cite{sommerville2011}:
\begin{description}
\item[a)] \textbf{Elicitação \& Análise}: A interface captura exatamente os insumos críticos para regressão espectral: \verb|Xcal|, \verb|Ycal|, \verb|Xteste|, \verb|Yteste| e os hiperparâmetros \textit{m} (bags) e \textit{k} (fração de subagging). O código reconhece peculiaridades do domínio, como suporte a decimal com vírgula (padrão BR) e detecção automática de separador CSV.
\item[b)] \textbf{Especificação}: Contratos de entrada/saída estão formalizados nas funções \verb|parse_matrix()|, \verb|parse_vector()|, \verb|perform_rlm()|, \verb|perform_bagging()| e \verb|calc_metrics()|. Todos operam com \verb|np.ndarray| ou tipos numéricos explícitos, garantindo precisão matemática e rastreabilidade.
\item[c)] \textbf{Validação}: A validação cruzada 5-fold (\verb|run_cross_validation|) e a comparação final de métricas (RMSE, MAE, R²) no conjunto de teste funcionam como \textit{oráculos} de equivalência. O log formatado \verb|result["log"]| documenta cada passo, permitindo auditoria científica direta e comparação com os resultados do sistema legado.
\end{description}

\section{Estratégias de Migração de Sistemas}
A transição do modelo Desktop (Java) para o modelo Web foi conduzida por reescrita funcional, focando na extração de requisitos sem reaproveitar o código-fonte antigo:
\begin{description}
\item[a)] \textbf{Reescrita Black-Box}: O comentário inicial \textit{``Conversão do projeto Vaadin Flow (Java) para Streamlit (Python)''} e as referências \textit{``espelho do RegressionService.java''} indicam uma migração por reimplementação. Não há dependência de JVM, nem reaproveitamento de \verb|.class| ou \verb|.jar|~\cite{becker2020}.
\item[b)] \textbf{Foco na Lógica de Negócio}: A migração extraiu apenas o núcleo matemático (pseudoinversa, amostragem com/sem reposição, métricas de erro) e o traduziu para \verb|numpy|, eliminando camadas de UI Java e substituindo-as por widgets nativos do Streamlit.
\item[c)] \textbf{Deploy Web Nativo}: O modelo Desktop/Servidor Java foi convertido para uma aplicação web \textit{stateless}, executável via \texttt{streamlit run app.py}, containerizável e acessível via navegador, reduzindo atrito de instalação e permitindo uso colaborativo em laboratórios.
\end{description}

\section{Mudança de Paradigma Arquitetural}
A arquitetura do sistema reflete uma mudança profunda no fluxo de execução e gerenciamento de estado:
\begin{description}
\item[a)] \textbf{MVC Tradicional (Java/Vaadin)}: Separação explícita entre \textit{Model} (dados/regressão), \textit{View} (componentes HTML/Java) e \textit{Controller} (listeners de evento). O estado vivia na sessão HTTP e as atualizações de interface eram parciais e assíncronas (AJAX/push).
\item[b)] \textbf{Streamlit (Fluxo de Dados Unidirecional)}: O \verb|app.py| é um script linear que \textbf{reexecuta do início ao fim} a cada interação. Não há controlador permanente; o fluxo segue \textit{Entrada do Usuário $\rightarrow$ Rerun do Script $\rightarrow$ Renderização Declarativa}. O estado é reconstruído e persistido apenas via \verb|st.session_state["result"]|, que armazena o dicionário de saída para que as abas \textbf{Resultados} e \textbf{Gráficos} leiam sem disparar novo cálculo.
\end{description}

\chapter{Ecossistema de Desenvolvimento (Python e Streamlit)}
\section{Python na Computação Científica}
O \verb|app.py| justifica a adoção do Python pelo ecossistema maduro voltado ao processamento de dados e álgebra linear:
\begin{description}
\item[a)] \textbf{Processamento Vetorial com \texttt{NumPy}}: Operações como \verb|np.linalg.pinv|, amostragem via \verb|np.random.default_rng| e cálculo de métricas (\verb|np.sqrt|, \verb|np.mean|) substituem laços explícitos, garantindo desempenho próximo a linguagens compiladas para matrizes espectrais.
\item[b)] \textbf{Parsing Robusto com \texttt{Pandas}}: As funções \verb|parse_matrix()| e \verb|parse_vector()| utilizam \verb|pd.read_csv| e \verb|pd.read_excel| com coercão automática (\path{pd.to_numeric(errors="coerce")}), limpando dados inválidos e normalizando formatos (\texttt{.txt}, \texttt{.csv}, \texttt{.xlsx}) sem perda de precisão.
\item[c)] \textbf{Interoperabilidade Científica}: A integração nativa com bibliotecas como \verb|plotly| e \verb|openpyxl| permite a transição direta de dados brutos para visualização e relatórios, mantendo o fluxo de trabalho típico de laboratórios e grupos de quimiometria.
\end{description}

\section{Ciclo de Vida do Streamlit}
O framework opera sob um modelo declarativo e \textit{stateless} por padrão, impactando diretamente a arquitetura do \verb|app.py|:
\begin{description}
\item[a)] \textbf{Execução Top-Down Unidirecional}: A cada interação do usuário (clique, alteração de input ou navegação entre abas), o interpretador relê o script do início ao fim. Não há um \textit{controller} permanente; o fluxo segue \textit{Entrada $\rightarrow$ Rerun $\rightarrow$ Renderização}.
\item[b)] \textbf{Persistência via \texttt{st.session\_state}}: Para evitar recálculos desnecessários, o dicionário \verb|st.session_state["result"]| armazena métricas, predições e o log completo. As abas \textbf{Resultados} e \textbf{Gráficos} leem esse estado sem acionar novamente a regressão.
\item[c)] \textbf{Feedback Assíncrono Simulado}: A barra de progresso (\verb|st.progress|) é atualizada via \texttt{callback} \verb|update_progress| dentro do laço de otimização. Embora o Streamlit seja síncrono por natureza, o \texttt{callback} mitiga a percepção de travamento durante a validação cruzada 5-fold.
\end{description}

\section{Gerenciamento de Memória e Performance}
A gestão eficiente de recursos no \verb|app.py| é crítica para manter a responsividade durante o processamento de matrizes espectrais:
\begin{description}
\item[a)] \textbf{Uso Estratégico de \texttt{st.session\_state}}: O estado mantém o output da análise vivo entre \textit{reruns}, eliminando a necessidade de reprocessar arquivos já carregados ou reexecutar a otimização de hiperparâmetros ao navegar pelas abas.
\item[b)] \textbf{Oportunidade de \texttt{@st.cache\_data}}: Atualmente, funções como \verb|parse_matrix()| e \verb|run_full_analysis()| são recalculadas a cada interação se o estado não for mantido. A aplicação de \verb|@st.cache_data| nessas funções reduziria a latência em 70--90\% para execuções repetidas com os mesmos inputs.
\item[c)] \textbf{Isolamento de Carga Computacional}: A lógica pesada (\verb|find_optimal_parameters|, \verb|perform_bagging|) está encapsulada em funções puras que recebem \verb|np.ndarray| e retornam resultados determinísticos. Isso evita que operações de UI bloqueiem o \textit{event loop} do Streamlit e facilita a futura migração para processamento em \textit{background} (\verb|concurrent.futures|).
\end{description}

\section{IDEs e Ambiente de Desenvolvimento}
O \verb|app.py| é desenvolvido e mantido com eficiência através do ecossistema do VS Code, padrão da indústria para projetos Python científicos:
\begin{description}
\item[a)] \textbf{Extensões Especializadas}: \texttt{Python}, \texttt{Pylance} (type checking para \verb|np.ndarray|), \texttt{Streamlit} (hot-reload e preview) e \texttt{Jupyter} (prototipagem de funções matemáticas antes da integração no script).
\item[b)] \textbf{Qualidade e Padronização}: Integração nativa com \texttt{black} (formatação), \texttt{ruff}/\texttt{flake8} (linting) e \texttt{mypy} (validação estática). O código já segue boas práticas (funções puras, docstrings, separação de camadas), reforçadas por diagnósticos em tempo real.
\item[c)] \textbf{Depuração e Ambientes Isolados}: O \texttt{launch.json} permite breakpoints dentro de \verb|run_cross_validation| ou \verb|calc_metrics|, inspecionando \textit{shapes} de arrays e valores de \verb|rmse_list| em tempo real. O gerenciamento de \texttt{venv}/\texttt{conda} via terminal integrado isola dependências (\verb|streamlit|, \verb|numpy|, \verb|plotly|) de conflitos do sistema operacional.
\end{description}

\chapter{Interface e Experiência do Usuário (UI/UX) Científica}
\section{Design Centrado em Dados}
O \verb|app.py| organiza o fluxo analítico em uma interface hierárquica, priorizando a clareza e a produtividade do cientista:
\begin{description}
\item[a)] \textbf{Navegação por Etapas (\texttt{st.tabs})}: O uso de abas (\textbf{Configuração}, \textbf{Resultados}, \textbf{Gráficos}) guia o usuário em um fluxo linear: entrada de dados $\rightarrow$ processamento $\rightarrow$ inspeção. Isso reduz a sobrecarga cognitiva e espelha o \textit{workflow} típico de validação de modelos quimiométricos.
\item[b)] \textbf{Layout em Colunas e Hierarquia Visual}: \verb|st.columns(2)| agrupa uploads de calibração e teste lado a lado. A área central é reservada para \textit{cards} de métricas e tabelas comparativas. O CSS customizado (\verb|.metric-card|, \verb|.log-area| estilo terminal) replica \textit{dashboards} técnicos e industriais.
\item[c)] \textbf{Divulgação Progressiva (\texttt{st.expander})}: Detalhes técnicos de formatação de arquivo (\verb|_EXP_MATRIZ|, \verb|_EXP_VETOR|) são ocultados por padrão e revelados sob demanda, mantendo a interface limpa enquanto garante acesso rápido à documentação de \textit{input}.
\end{description}

\section{Visualização Interativa com Plotly}
A escolha do Plotly no \verb|app.py| é fundamentada pela necessidade de inspeção técnica precisa em espectrometria e regressão:
\begin{description}
\item[a)] \textbf{Integração Nativa e Otimizada}: \texttt{st.plotly\_chart(..., use\_container\_width=True)}
 renderiza gráficos vetoriais diretamente no navegador, sem conversão adicional ou perda de qualidade, mantendo a responsividade mesmo com milhares de pontos.
\item[b)] \textbf{Ferramentas de Inspeção Científica}: Os gráficos \verb|go.Scatter| habilitam zoom, \textit{pan}, \textit{hover} com coordenadas exatas e legendas interativas. A linha ideal \textit{y=x} (\texttt{mode="lines", dash="dot"}) no gráfico de dispersão permite validação visual imediata de viés e dispersão dos modelos.
\item[c)] \textbf{Padrões de Publicação}: Eixos customizados (\verb|showgrid=True|, \verb|gridcolor="#000000"|, fontes pretas explícitas) garantem legibilidade para relatórios técnicos, apresentações ou submissão a periódicos, alinhando-se às normas de visualização de dados científicos.
\end{description}

\section{Contratos de Interface e Funções}
O \verb|app.py| aplica rigorosamente a separação de responsabilidades, assegurando robustez, testabilidade e evolução contínua:
\begin{description}
\item[a)] \textbf{Separação de Camadas}: As funções matemáticas (\verb|perform_rlm|, \verb|perform_bagging|, \verb|calc_metrics|, \verb|run_cross_validation|) operam exclusivamente com \verb|np.ndarray| e tipos numéricos. Elas \textbf{não importam} \verb|streamlit|, podendo ser testadas isoladamente com \texttt{pytest} e dados sintéticos. A camada de UI cuida apenas de estado, validação de \textit{input} e renderização.
\item[b)] \textbf{Contratos Explícitos e Validação}: Funções como \verb|parse_int_list()| e \verb|parse_float_list()| lançam \verb|ValueError| com mensagens claras para entradas inválidas. A tipagem estática e a coercão automática (\verb|pd.to_numeric|) previnem falhas silenciosas durante o parse de arquivos brutos.
\item[c)] \textbf{Manutenibilidade e Escalabilidade}: Alterar um algoritmo de otimização, adicionar uma nova métrica ou modificar o layout CSS não impacta as outras camadas. A arquitetura modular permite a futura extração da lógica de regressão para um pacote Python independente, enquanto o \verb|app.py| atua apenas como \textit{frontend} web.
\end{description}